%% ------------- English version ---------------
\documentclass[english,hidelinks]{sbrt}
\usepackage{newtxtext}
\usepackage[english]{babel}
\usepackage[utf8]{inputenc}
\newtheorem{theorem}{Theorem}
%% ---------------------------------------------
\setlength{\marginparwidth}{2cm}
\usepackage{todonotes}
\newcommand\rigel[1]{\todo[color=green, inline]{\textbf{Rigel}: #1}}
\usepackage{orcidlink}
\usepackage{hyperref}
\usepackage[
  compatibility=false,
  font=footnotesize,
  labelsep=period,
  format=plain,
  justification=raggedright,
  singlelinecheck=false
]{caption}

\begin{document}

\title{Low-Cost Embedded System for Beehive Monitoring}

\author{Igor Costa, Jorge Alves, Ian Dias and Davi Ito
  \thanks{The authors are with Ibmec, Rio de Janeiro, RJ, Brazil.
  E-mails: igorcosta.silva10a@gmail.com,
  jorgedocruzeiro07@gmail.com,
  iandpcarvalho@gmail.com,
  cdavi8249@gmail.com.}%
  }

\maketitle

\markboth{XLIV BRAZILIAN SYMPOSIUM ON TELECOMMUNICATIONS AND SIGNAL PROCESSING - SBrT 2026, SEPTEMBER 29TH TO OCTOBER 2ND, 2026, SALVADOR, BA}{}

\begin{abstract}
\end{abstract}
\begin{keywords}
Beehive monitoring, embedded systems, Internet of Things, computer vision.
\end{keywords}

\section{Introduction}

Bees play a central role in natural and agricultural ecosystems because pollination by bees supports the reproduction of flowering plants and the production of crops. The western honey bee (\textit{Apis mellifera}) is among the most widely distributed and frequently observed crop pollinators, while its importance in natural habitats can vary according to the ecological context. At the same time, bee populations face multiple pressures, including habitat loss, invasive species, emerging diseases, pesticide exposure, and climate change. These conditions strengthen the need for reliable observation of colony and bee activity, especially when changes must be identified over time rather than through isolated inspections \cite{hung2018pollinators,brown2009conservation}.

Traditional hive inspections provide valuable information about colony conditions, but they are intermittent, require human presence, and may disturb the colony when the hive is opened repeatedly. Embedded systems are computing platforms designed for specific functions and commonly optimized for constraints such as processing time, reliability, power consumption, size, and cost. Internet of Things (IoT) architectures extend these systems by connecting sensors and other physical devices to communication networks and data services, allowing data to be processed locally, remotely, or across both levels. De Micco et al. reviewed the characteristics and applications of embedded systems, including their relation to IoT and edge computing, while Ray surveyed IoT architectures composed of devices, sensors, communication protocols, services, and processing layers. This technological foundation motivates a low-cost, non-invasive beehive monitoring system in which environmental and visual measurements are collected by embedded devices and made available for subsequent analysis \cite{demicco2020embedded,ray2018iotarchitectures}.

\textit{Related work.} Entrance-focused computer vision has been combined with environmental sensing to measure colony activity without opening the hive. Zheng et al. placed temperature and humidity sensors at the hive center, margin, entrance, and surrounding environment and used an internal sound sensor and an entrance-facing webcam. Their YOLOv5--DeepSORT system reached 83.5\% $\pm$ 0.7 MOTA and 77.3\% $\pm$ 0.2 MOTP at up to 16 frames per second. Narcia-Macias et al. reported 96.28\% tracking accuracy and an F1-score of 0.8319 for pollen detection with the low-cost YOLOv7-tiny-based IntelliBeeHive. However, the mite data were insufficient and partly based on placeholder data, so Varroa detection should be treated as a proof of concept rather than a validated field capability. Together, these studies support entrance imaging and lightweight tracking while showing that each target task requires representative and correctly labeled data \cite{zheng2024intelligent,narciamacias2024intellibeehive}.

Acoustic and vibration sensing provide another non-invasive approach for assessing colony conditions without repeatedly opening the hive. Uthoff et al. reviewed the use of microphones and accelerometers to investigate beekeeping-relevant states associated with queen presence and swarming preparation. The review identifies promising relationships between colony-generated signals and these states, but also highlights small datasets, the absence of standardized evaluation metrics, inconsistent feature-engineering procedures, and limited generalizability across colonies and recording conditions. These limitations indicate that acoustic or vibration indicators in the proposed project should initially be treated as experimental evidence and validated under the sensors, hive geometry, and environmental conditions available to the group \cite{uthoff2023acoustic}.

IoT platforms and wireless sensor networks support continuous acquisition and remote access to measurements collected inside the hive. Tashakkori et al. presented Beemon, an architecture that combines temperature, humidity, weight, audio, and video acquisition with MQTT communication, a ThingsBoard dashboard, and remote storage. Henry et al. developed a wireless sensor network for continuous in-hive temperature, relative-humidity, and acoustic measurements in an apiary environment. These systems demonstrate the feasibility of integrating multiple sensing modalities and maintaining a historical data stream, while also motivating the proposed project to document timestamps, communication reliability, sensor placement, and missing measurements instead of treating data transmission alone as evidence of colony health \cite{tashakkori2021beemon,henry2019precision}.

Tu et al. investigated lightweight entrance monitoring through computer vision implemented on a Raspberry Pi. Their background-subtraction approach counted bees, identified their positions, and estimated entrance and exit activity, reporting $R^2$ values of 0.987 for bee counting, 0.953 for incoming activity, and 0.888 for outgoing activity. The results support local processing on a low-cost platform and the use of manually annotated counts as an evaluation reference. Nevertheless, performance may vary with illumination, camera placement, entrance geometry, and overlapping bees. For the proposed system, this work motivates scheduled entrance-image acquisition and reproducible comparison with manual observations before more complex detection or tracking models are adopted \cite{tu2016automatic}.

The implementation, documentation, and experimental material produced in this project are publicly available in the project repository \cite{repo}.

This work proposes a low-cost, non-invasive monitoring station that records entrance activity and environmental conditions of a beehive without requiring any network infrastructure at the apiary. A camera mounted above and slightly in front of the entrance captures short image sequences on a fixed schedule, two temperature and humidity sensors record the interior and the exterior of the hive, and both streams are written with synchronised timestamps to a memory card carried by the station, which is powered by a small solar panel and a rechargeable battery. The recorded material is collected weekly and analysed offline, where a detection model counts the bees visible in each observation window. The contributions of this work are fourfold. First, we describe a station that operates autonomously and stores its data locally, so that data collection does not depend on wireless coverage, which is frequently unavailable where hives are kept; this is a deliberate departure from the networked architectures of Tashakkori et al. and Henry et al. \cite{tashakkori2021beemon,henry2019precision}. Second, we define an acquisition protocol in which two capture configurations are compared empirically by counting quality, storage footprint, and energy consumption, so that the sampling interval is reported as an engineering trade-off rather than as an arbitrary setting. Third, we pair every observation window with interior and exterior temperature and humidity, which allows the observed activity to be read against the conditions under which it was recorded. Fourth, we evaluate the counts against manual annotations of the same images, following the reproducible comparison adopted by Tu et al. for low-cost entrance monitoring \cite{tu2016automatic}. The analysis is deliberately restricted to counting the bees visible in each window as an indicator of entrance activity. Because the camera is inactive between windows, the system does not attempt to follow an individual bee across windows, and it makes no claim about pollen transport, \textit{Varroa} infestation, queen presence, or swarming preparation, all of which require sensing modalities and labelled data that the first version does not provide.

The rest of this work is organized as follows. Section~\ref{sec:system-description} describes the monitoring station, presenting the camera, the environmental sensors, the local storage and timekeeping, and the solar power supply. Section~\ref{sec:method} details the proposed method, covering the acquisition schedule, the two capture configurations under comparison, and the offline counting pipeline. Section~\ref{sec:results} presents the qualitative and quantitative results obtained with the prototype. Section~\ref{sec:conclusion} concludes the work and outlines the next steps.

\rigel{Pendência da AC Semana 5: verificar as referências produzidas por alunos do Ibmec e citá-las caso seja oportuno.}

\section{System Description}
\label{sec:system-description}

\rigel{AC Semana 6: Apresente os sensores e/ou microcontroladores e, possivelmente, os protocolos de comunicação utilizados. Inclua uma referência que apresente cada um desses componentes (pode ser referências já utilizadas previamente na seção de introdução) que deverão dar suporte a essa descrição.}

\subsection{Local Storage and Offline Operation}

The proposed system does not depend on a network connection to operate. Environmental measurements and entrance images are written to a microSD card carried by the acquisition board, and a real-time clock keeps the timestamps consistent even when no external time source is available. The card is retrieved approximately once a week, and the recorded material is analysed on a separate computer.

This decision follows from the deployment context. Apiaries are frequently located where wireless coverage is unreliable or absent, and a monitoring station that stops recording when the link fails would produce exactly the gaps that continuous observation is meant to eliminate. Tashakkori et al. and Henry et al. demonstrate the benefits of networked acquisition through MQTT and wireless sensor networks, and remote transmission remains a natural extension of this work; for the first version, however, local storage removes the communication subsystem from the set of components that can interrupt data collection \cite{tashakkori2021beemon,henry2019precision}.

\section{The Method}
\label{sec:method}

\rigel{AC Semana 7:
Nesse momento já teremos a seção Introdução, a seção II e a Seção III que trata do método que o grupo propõe.
Atividade relativa a parte escrita: Elaborar a seção III que deverá tratar sobre o método proposto:

(1) Apresente, em detalhes, o funcionamento do sistema desenvolvido. Como nessa seção estão as contribuições do grupo, não é necessário inserir referências que dêem suporte;

(2) Elabore um diagrama que facilite a compreensão do método (use o LucidChart por exemplo, exporte o diagrama como PDF para manter a boa resolução no artigo, cuidado com os espaços em branco que podem sobrar às margens do diagrama (no momento de exportar como PDF faça um recorte personalizado para tirar as margens em branco indesejáveis).

De forma alternativa (ou adicionalmente), o grupo poderá gerar um pseudocódigo para mostrar como o sistema funciona (passos principais do método proposto);

(3) Um parágrafo deverá ser escrito e deverá referenciar o diagrama criado na seção III Método. Explique o diagrama criado ao leitor.
(4) Outro parágrafo deverá ser escrito e referenciar uma foto do protótipo. Explique o que está na foto ao leitor.}

\section{Results}
\label{sec:results}

\rigel{AC Semana 8:
Nesse momento já teremos a seção Introdução, a seção II e a Seção III que trata do método que o grupo propõe.
Atividade relativa a parte escrita: Elaborar a seção IV que deverá tratar sobre os resultados experimentais:

(1) Apresente, em detalhes, os resultados qualitativos e quantitativos obtidos até o momento. No mínimo resultados qualitativos.

(2) Elabore o Abstract e Conclusions.}

\rigel{AC Semana 9: Continuar na elaboração da seção IV que deverá tratar sobre os resultados experimentais:

(1) Introduza na Seção IV (Results), em detalhes, os resultados quantitativos obtidos (além de resultados qualitativos já conseguidos).

(2) Refine o abstract. Na última linha do abstract devemos inserir algum resultado que seja interessante para o leitor e que represente algum aspecto considerado como a maior contribuição do grupo

(3) Refine a seção Conclusions. Insira resultados que resumam a avaliação do sistema embarcado proposto.}

\section{Conclusion}
\label{sec:conclusion}

\bibliographystyle{ieeetr}
\bibliography{refs}

\end{document}